\documentclass[../DoAn.tex]{subfiles}
\begin{document}
Phụ lục này trình bày thiết lập và kết quả chi tiết của các bộ đánh giá tự động được tóm tắt ở Chương \ref{chapter:SolutionAndContribution}. Mã nguồn các bộ đánh giá nằm trong thư mục \texttt{evaluation} và kết quả gốc dạng JSON nằm trong \texttt{evaluation/\allowbreak{}results}.

\section{Đánh giá truy xuất}
\textbf{Thiết lập.} Bộ đánh giá \texttt{eval\_retrieval\_nogt.py} chạy mỗi cấu hình truy xuất trên 50 câu hỏi và lấy 10 kết quả đầu. Kết quả của mọi cấu hình được gộp và khử trùng; thứ tự các cặp được xáo trộn với hạt giống 7 và mỗi cặp câu hỏi và điều luật được một LLM (Claude Sonnet, mức suy luận thấp) chấm mù theo thang 2 (trả lời trực tiếp câu hỏi), 1 (liên quan hoặc hữu ích một phần), 0 (không liên quan), không biết cặp đến từ cấu hình nào. Kết quả chấm được lưu đệm để chạy lại không tốn thêm chi phí. P@5, P@10 và Top1 tính với mức liên quan từ 1 trở lên; MRR@10 tính với kết quả có mức liên quan bằng 2 đầu tiên; nDCG@10 dùng thứ tự lý tưởng là các điều tốt nhất trong toàn bộ kết quả gộp của lần chạy, nên chỉ so được giữa các cấu hình trong cùng một lần chạy. Các cấu hình được so theo từng câu bằng nDCG@10 và kiểm định dấu hai phía.

Bảng \ref{tab:pl_sign} so sánh từng câu với cấu hình ban đầu; Bảng \ref{tab:pl_sign_hybrid} so sánh các biến thể với cấu hình lai trong cùng lần chạy.

\input{tables/tab_sign_old}

\input{tables/tab_sign_hybrid}

\textbf{Đánh giá theo mục đã biết.} Bộ đánh giá chọn ngẫu nhiên 100 điều luật có nội dung dài hơn 200 ký tự (hạt giống 42), dùng GPT-4o-mini sinh một câu hỏi đời thường cho mỗi điều và coi chính điều đó là đáp án. Câu hỏi được lưu trong \texttt{evaluation/\allowbreak{}synthetic\_queries.json} nên các lần chạy dùng cùng một bộ. Vì câu hỏi được sinh từ nội dung điều, phép đo này thiên về BM25. Kết quả được trình bày cùng Bảng \ref{tab:kq_truyxuat}.

\textbf{Thử nghiệm ban đầu.} Trước khi xây dựng bộ đánh giá trên, đồ án thử nghiệm bốn chiến lược trên năm câu hỏi có đáp án gán tay, đo bằng P@10, Hit@10, nDCG@10 và MRR (Bảng \ref{tab:dot1}). Với năm câu hỏi, kết quả chỉ được dùng để định hướng; đồ án không dùng các chỉ số so khớp văn bản như ROUGE-L hay BLEU để kết luận về chất lượng.

\input{tables/tab_dot1}

\section{Đánh giá hành vi gọi công cụ}
Bộ đánh giá \texttt{eval\_tool\_calls.py} chạy vòng lặp tác tử thật trên các tình huống ghi trong \texttt{tool\_call\_cases.json}. Mỗi tình huống khai báo công cụ bắt buộc phải gọi, công cụ nên gọi, công cụ không được gọi, công cụ phải gọi đầu tiên, việc luật buộc web có được kích hoạt hay không, trần số lần gọi và các từ khoá câu trả lời phải chứa. Chỉ số ``mã bịa đặt'' đo tỷ lệ lượt gọi đồ thị hoặc đọc toàn văn với mã điều chưa từng xuất hiện trong kết quả trước đó. Bảng \ref{tab:pl_tool} trình bày kết quả của hai mô hình tác tử trên cùng mã nguồn; mỗi tình huống được chạy một lần.

\input{tables/tab_toolcall}

Theo nhóm câu hỏi, số tình huống đạt với GPT-4.1 lần lượt là: định nghĩa 5/5, phạm vi áp dụng 4/5, mức phạt 6/6, ngoài dữ liệu 2/3, quan hệ giữa điều luật 1/1, đọc toàn văn 1/1, trò chuyện 2/2. GPT-4o-mini đạt như GPT-4.1 ở mọi nhóm trừ nhóm ngoài dữ liệu (0/3).

\section{Đánh giá bộ nhớ dài hạn}
Bộ đánh giá \texttt{eval\_memory.py} tạo người dùng giả riêng cho từng kịch bản, chạy chuỗi lượt hội thoại qua mô hình và MongoDB thật, kiểm tra kỳ vọng rồi xoá dữ liệu. Các kỳ vọng gồm: bộ nhớ phải chứa hoặc không được chứa các cụm từ nhất định, thông tin phải thuộc đúng loại, tổng số ký ức tối đa, số ký ức trùng khớp tối đa (phát hiện bản trùng chưa xoá) và truy vấn phải tìm lại được thông tin cần thiết. Bảng \ref{tab:pl_mem} liệt kê 14 kịch bản và Bảng \ref{tab:pl_mem_runs} trình bày kết quả các lần chạy.

\begin{table}[H]
\centering
\caption{Các kịch bản đánh giá bộ nhớ dài hạn}
\label{tab:pl_mem}
\begin{tabular}{|L{4.6cm}|L{2.2cm}|L{6.6cm}|}
\hline
\textbf{Kịch bản} & \textbf{Nhóm} & \textbf{Kỳ vọng chính} \\ \hline
\texttt{profile\_basic} & Trích xuất & Lưu năm sinh, nghề nghiệp, nơi ở \\ \hline
\texttt{chitchat\_nothing} & Trích xuất & Trò chuyện xã giao, không lưu gì \\ \hline
\texttt{general\_law\_not\_profile} & Trích xuất & Câu hỏi luật chung không bị coi là thông tin cá nhân \\ \hline
\texttt{multi\_fact\_turn} & Trích xuất & Một lượt chứa nhiều thông tin, tách đủ \\ \hline
\texttt{episodic\_wage} & Trích xuất & Sự kiện bị nợ lương lưu đúng loại sự kiện \\ \hline
\texttt{procedural\_pref} & Trích xuất & Mong muốn cách trả lời lưu đúng loại thói quen \\ \hline
\texttt{sensitive\_id\_number} & Nhạy cảm & Không lưu số căn cước \\ \hline
\texttt{sensitive\_bank} & Nhạy cảm & Không lưu số tài khoản ngân hàng \\ \hline
\texttt{update\_school} & Cập nhật & Đổi trường thì cập nhật, không giữ bản cũ \\ \hline
\texttt{update\_job} & Cập nhật & Đổi việc thì cập nhật \\ \hline
\texttt{negation\_delete} & Cập nhật & Người dùng phủ định thông tin cũ thì xoá \\ \hline
\texttt{paraphrase\_noop} & Cập nhật & Nói lại cùng thông tin thì bỏ qua \\ \hline
\texttt{retrieve\_age\_for\_marriage} & Truy xuất & Tìm lại năm sinh khi hỏi điều kiện kết hôn \\ \hline
\texttt{retrieve\_relevant\_situation} & Truy xuất & Tìm lại tình huống liên quan \\ \hline
\end{tabular}
\end{table}

\input{tables/tab_memory}

\section{Đánh giá kiểm tra ngoại lệ}
Bộ đánh giá \texttt{eval\_exceptions.py} sinh các cặp điều chính và điều ngoại lệ từ đồ thị như mô tả ở phần \ref{section:5.graph}; 30 cặp đã sinh được lưu trong \texttt{evaluation/\allowbreak{}exception\_cases.json} nên các lần chạy sau không gọi lại Claude. Ở chế độ truy xuất, mỗi câu hỏi được viết lại bằng GPT-4o-mini như trong hệ thống rồi đưa vào \texttt{vector\_search} với kiểm tra ngoại lệ tắt và bật; một cặp được tính là có điều A hoặc B nếu mã điều tương ứng xuất hiện trong kết quả của công cụ. Kết quả ở Bảng \ref{tab:ngoaile}.

\end{document}
